\documentclass[../main.tex]{subfiles}

\begin{document}

\section{Introduction}\label{sec:introduction}

This report presents the solution to Homework~1 of the course \emph{Dynamics and Control of Launch Vehicles}, which concerns the minimum-fuel planar ascent of a rocket from rest on the ground to a prescribed low-Earth orbit injection state. Four progressively richer mission programs are addressed, each framed as an optimal control problem and solved by applying the Pontryagin Maximum Principle (PMP) and single-shooting over the costate initial conditions.

\subsection{Indirect Optimal Control: Method Overview}\label{sec:method}

All four mission programs are solved by the \emph{indirect} approach to optimal control, which follows an \emph{optimize-then-discretize} paradigm: the first-order necessary conditions of optimality are derived analytically from the Pontryagin Maximum Principle~\cite{pontryagin1962} and only the resulting boundary value problem is solved numerically. Concretely, the workflow is a four-step pipeline:
\begin{enumerate}
    \item form the Hamiltonian and impose its pointwise maximization to express the optimal control as a function of state and costate;
    \item derive the costate (adjoint) dynamics $\dot{\boldsymbol{\lambda}} = -\partial\mathcal{H}/\partial\mathbf{x}$ and the transversality conditions that supply the missing boundary data;
    \item assemble the resulting two-point boundary value problem (BVP) for the augmented state--costate system;
    \item solve the BVP by single shooting over the unknown initial costates and the free final time.
\end{enumerate}
The competing \emph{direct} approach reverses this order---it discretizes the trajectory first and optimizes second, transcribing it into a large nonlinear program~\cite[Sec.~4.12.7]{betts2010}. The indirect route adopted here is attractive on three counts: it is highly accurate, since the optimality conditions are imposed exactly rather than approximated on a mesh; it carries a transparent physical structure, the velocity costate $\boldsymbol{\lambda}_v$ being the \emph{primer vector}~\cite{lawden1963} that fixes the optimal thrust direction; and it collapses the problem to a handful of unknowns. Its price is the notorious sensitivity of the shooting residual to the costate guess---which has no direct physical interpretation---together with small convergence basins and the need to postulate the thrust-arc (switching) structure a priori. Section~\ref{sec:shooting} returns to these difficulties and to the continuation strategy used to overcome them.

\subsection{Vehicle Dynamics}\label{sec:dynamics}

The vehicle is modelled as a point mass moving in a vertical plane above a flat, non-rotating Earth with no atmospheric drag. The state vector is
\begin{equation}
    \mathbf{x}(t) = \bigl[x,\; y,\; v_x,\; v_y,\; m\bigr]^T,
\end{equation}
where $(x, y)$ is the position in the orbital plane, $(v_x, v_y)$ is the velocity, and $m$ is the vehicle mass. The single control variable is the thrust-direction angle $\varphi$ measured from the horizontal. The thrust magnitude is $T = cQ$, where $c$ is the effective exhaust velocity and $Q = -\dot{m}$ is the constant mass-flow rate. The equations of motion are
\begin{align}
    \dot{x}   &= v_x, \label{eq:xdot}\\
    \dot{y}   &= v_y, \label{eq:ydot}\\
    \dot{v}_x &= \frac{T}{m}\cos\varphi, \label{eq:vxdot}\\
    \dot{v}_y &= \frac{T}{m}\sin\varphi - g, \label{eq:vydot}\\
    \dot{m}   &= -Q. \label{eq:mdot}
\end{align}

\subsection{Non-dimensionalization}\label{sec:nondim}

All quantities are made dimensionless with the reference values in Table~\ref{tab:nondim}. The reference velocity $V_\mathrm{ref}$ is the circular orbital speed at the nominal target altitude; the reference length and time follow from $g$.

\begin{table}[H]
\centering
\caption{Non-dimensionalization reference values.}
\label{tab:nondim}
\begin{tabular}{llll}
\toprule
Quantity     & Symbol              & Expression              & Value \\
\midrule
Velocity     & $V_\mathrm{ref}$   & orbital speed           & 7800 m/s \\
Acceleration & $g_\mathrm{ref}$   & $g$                     & 9.81 m/s$^2$ \\
Length       & $L_\mathrm{ref}$   & $V_\mathrm{ref}^2/g$    & $\approx 6200$ km \\
Time         & $t_\mathrm{ref}$   & $V_\mathrm{ref}/g$      & $\approx 795$ s \\
Mass         & $m_\mathrm{ref}$   & $m_0$                   & $10^6$ kg \\
\bottomrule
\end{tabular}
\end{table}

In non-dimensional form ($g = 1$, $m_0 = 1$, $v_{x,f} = 1$) the fixed problem data are collected in Table~\ref{tab:params}.

\begin{table}[H]
\centering
\caption{Fixed non-dimensional problem parameters.}
\label{tab:params}
\begin{tabular}{lll}
\toprule
Symbol   & Value             & Physical interpretation \\
\midrule
$c$      & 0.6               & Effective exhaust velocity ($c_\mathrm{dim} \approx 4680$ m/s, $I_{sp} \approx 477$ s) \\
$\eta$   & 0.1               & Structural coefficient $m_s/m_p$ \\
$y_f$    & 0.04, 0.05, 0.06  & Target orbit altitude ($\approx 248$, 310, 372 km) \\
\bottomrule
\end{tabular}
\end{table}

The mass-flow rate $Q$ is a free design parameter swept in Tasks~1 and~4. Since $W = mg$ and $m_0 = 1$, the initial thrust-to-weight ratio is $\left.T/W\right|_{t=0} = cQ$, so liftoff ($T > W$) requires $Q > 1/c \approx 1.67$. As propellant burns off ($\dot{m} = -Q < 0$) the ratio $T/W = cQ/m$ increases monotonically, hence the liftoff instant is the binding case.

The value $Q = 2$ is the nominal mass-flow rate \emph{given in the assignment data}. It corresponds to a liftoff thrust-to-weight ratio $\left.T/W\right|_{t=0} = cQ = 1.2$---representative of a real first stage and safely above the $Q_{\min} = 1/c \approx 1.67$ liftoff threshold. Task~1 sweeps $Q$ to study its effect and finds the payload-optimal rate $Q^\ast \approx 2.2$--$2.5$ (Table~\ref{tab:qstar}), so the nominal value sits just below the single-stage optimum. Tasks~2--4 then keep $Q = 2$ fixed, so that the comparison of the structural features (vertical climb, coast, staging) is made at a single, common operating point.

\subsection{Pontryagin Maximum Principle}\label{sec:pmp}

\subsubsection{Hamiltonian and Costate Equations}

The optimal control problem is to find the control $\varphi(t)$ and the free terminal time $t_f$ that maximize the final mass $m(t_f)$---equivalently, minimize the Mayer cost $J = -m(t_f)$---subject to the dynamics \eqref{eq:xdot}--\eqref{eq:mdot} and the terminal constraints
\begin{equation}
    y(t_f) = y_f, \qquad v_x(t_f) = 1, \qquad v_y(t_f) = 0.
    \label{eq:terminal}
\end{equation}
The horizontal position $x(t_f)$ is free, and $t_f$ is free.

Introducing the costate (adjoint) vector $\boldsymbol{\lambda} = [\lambda_x,\; \lambda_y,\; \lambda_{v_x},\; \lambda_{v_y},\; \lambda_m]^T$, the Hamiltonian is
\begin{equation}
    \mathcal{H} = \lambda_x v_x + \lambda_y v_y
    + \lambda_{v_x}\frac{T}{m}\cos\varphi
    + \lambda_{v_y}\!\left(\frac{T}{m}\sin\varphi - 1\right)
    - \lambda_m Q.
    \label{eq:hamiltonian}
\end{equation}
Rearranging,
\begin{equation}
    \mathcal{H} = \underbrace{\lambda_x v_x + \lambda_y v_y - \lambda_{v_y} - \lambda_m Q}_{\text{state-only part}}
    + \frac{T}{m}\underbrace{\bigl(\lambda_{v_x}\cos\varphi + \lambda_{v_y}\sin\varphi\bigr)}_{\text{control-dependent part}}.
\end{equation}

The Euler--Lagrange costate equations $\dot{\boldsymbol{\lambda}} = -\partial\mathcal{H}/\partial\mathbf{x}$~\cite[Ch.~2]{bryson1975} yield
\begin{align}
    \dot{\lambda}_x     &= 0
        &&\Longrightarrow\quad \lambda_x = 0
        \quad\text{(}x(t_f)\text{ free } \Rightarrow \lambda_x(t_f)=0\text{)},
        \label{eq:lx}\\
    \dot{\lambda}_y     &= 0
        &&\Longrightarrow\quad \lambda_y = \mathrm{const},
        \label{eq:ly}\\
    \dot{\lambda}_{v_x} &= -\lambda_x = 0
        &&\Longrightarrow\quad \lambda_{v_x} = \lambda_{v_x,0} = \mathrm{const},
        \label{eq:lvx}\\
    \dot{\lambda}_{v_y} &= -\lambda_y
        &&\Longrightarrow\quad \lambda_{v_y}(t) = \lambda_{v_y,0} - \lambda_y t,
        \label{eq:lvy}\\
    \dot{\lambda}_m     &= \frac{T}{m^2}\left\|\boldsymbol{\lambda}_v\right\|,
        \label{eq:lm}
\end{align}
where $\|\boldsymbol{\lambda}_v\| = \sqrt{\lambda_{v_x}^2 + \lambda_{v_y}^2}$.

Four of the five costate components are determined analytically: $\lambda_x = 0$ identically, $\lambda_y$ and $\lambda_{v_x}$ are constant, and $\lambda_{v_y}$ varies linearly in time. Only $\lambda_m$ requires numerical integration via Eq.~\eqref{eq:lm}. This analytical structure greatly reduces the cost of evaluating the shooting residual.

\subsubsection{Optimal Control Law}

By the PMP, the optimal control $\varphi^*$ maximizes $\mathcal{H}$ pointwise. Since $T/m > 0$, this is equivalent to maximizing the control-dependent part of Eq.~\eqref{eq:hamiltonian},
\begin{equation}
    g(\varphi) = \lambda_{v_x}\cos\varphi + \lambda_{v_y}\sin\varphi,
    \label{eq:gphi}
\end{equation}
over $\varphi$. The first-order stationarity condition reads
\begin{equation}
    \frac{\partial\mathcal{H}}{\partial\varphi} =
    \frac{T}{m}\bigl(-\lambda_{v_x}\sin\varphi + \lambda_{v_y}\cos\varphi\bigr) = 0
    \qquad\Longrightarrow\qquad
    \tan\varphi^* = \frac{\lambda_{v_y}(t)}{\lambda_{v_x}}.
    \label{eq:bilinear}
\end{equation}
This is the \textbf{bilinear-tangent (linear-tangent) law}~\cite[Sec.~2.4]{bryson1975}: since $\lambda_{v_y}$ is linear in $t$ (Eq.~\ref{eq:lvy}) and $\lambda_{v_x}$ is constant, $\tan\varphi^*$ varies linearly with time.

\textbf{Geometric interpretation and branch selection.}\ Equation~\eqref{eq:bilinear} determines $\varphi^*$ only up to an additive $\pi$: it admits two roots per period. Geometrically, $\lambda_{v_x}$ and $\lambda_{v_y}$ are the two legs of a right triangle whose hypotenuse is the norm $\|\boldsymbol{\lambda}_v\| = \sqrt{\lambda_{v_x}^2 + \lambda_{v_y}^2}$, so the two candidate directions are
\begin{equation}
    \text{(I)}\;\;
    \cos\varphi^* = \frac{\lambda_{v_x}}{\|\boldsymbol{\lambda}_v\|},\;
    \sin\varphi^* = \frac{\lambda_{v_y}}{\|\boldsymbol{\lambda}_v\|}
    \qquad\text{or}\qquad
    \text{(II)}\;\;
    \cos\varphi^* = -\frac{\lambda_{v_x}}{\|\boldsymbol{\lambda}_v\|},\;
    \sin\varphi^* = -\frac{\lambda_{v_y}}{\|\boldsymbol{\lambda}_v\|}.
    \label{eq:two_cases}
\end{equation}
Substituting each branch into Eq.~\eqref{eq:gphi} discriminates between them:
\begin{equation}
    g_\mathrm{I} = \frac{\lambda_{v_x}^2 + \lambda_{v_y}^2}{\|\boldsymbol{\lambda}_v\|}
    = \frac{\|\boldsymbol{\lambda}_v\|^2}{\|\boldsymbol{\lambda}_v\|} = +\|\boldsymbol{\lambda}_v\|,
    \qquad
    g_\mathrm{II} = -\|\boldsymbol{\lambda}_v\|.
\end{equation}
Branch~(I) yields the maximum and branch~(II) the minimum, so the PMP selects branch~(I): the optimal thrust direction is aligned with the velocity costate $\boldsymbol{\lambda}_v$ (primer vector), and the control-dependent term attains its largest value,
\begin{equation}
    \lambda_{v_x}\cos\varphi^* + \lambda_{v_y}\sin\varphi^* = \|\boldsymbol{\lambda}_v\|.
    \label{eq:maxcond}
\end{equation}
Using this result, the optimal Hamiltonian simplifies to
\begin{equation}
    \mathcal{H}^* = \lambda_y v_y + \frac{T}{m}\|\boldsymbol{\lambda}_v\| - \lambda_{v_y} - \lambda_m Q.
    \label{eq:ham_opt}
\end{equation}

\subsubsection{Transversality Conditions}\label{sec:transversality}

The costate boundary conditions follow from the transversality (natural boundary) conditions of the PMP~\cite[Ch.~2]{bryson1975}. Let $\mathcal{M} = m(t_f)$ denote the terminal (Mayer) payoff to be maximized, and adjoin the terminal constraints $\boldsymbol{\chi}(\mathbf{x}(t_f)) = [\,y - y_f,\; v_x - 1,\; v_y\,]^T = \mathbf{0}$ with multipliers $\boldsymbol{\mu}$. The augmented payoff is
\begin{equation}
    \bar{\mathcal{M}} = \mathcal{M} + \boldsymbol{\mu}^T\boldsymbol{\chi}\bigl(\mathbf{x}(t_f)\bigr)
    + \int_{t_0}^{t_f}\!\boldsymbol{\lambda}^T\!\bigl(\mathbf{f} - \dot{\mathbf{x}}\bigr)\,\mathrm{d}t,
    \label{eq:aug_payoff}
\end{equation}
whose integral vanishes on any feasible trajectory ($\dot{\mathbf{x}} = \mathbf{f}$) and therefore leaves the optimum unchanged. Integrating the $\boldsymbol{\lambda}^T\dot{\mathbf{x}}$ contribution by parts and collecting the terms evaluated at $t_f$, the first variation of $\bar{\mathcal{M}}$ accumulates the terminal boundary term
\begin{equation}
    \delta\bar{\mathcal{M}}\big|_{t_f} =
    \left[\frac{\partial \mathcal{M}}{\partial \mathbf{x}_f}
    + \left(\frac{\partial \boldsymbol{\chi}}{\partial \mathbf{x}_f}\right)^{\!T}\!\boldsymbol{\mu}
    - \boldsymbol{\lambda}(t_f)\right]^{\!T}\!\delta\mathbf{x}_f ,
\end{equation}
which must vanish for arbitrary admissible $\delta\mathbf{x}_f$. This yields the transversality condition
\begin{equation}
    \boldsymbol{\lambda}(t_f) = \frac{\partial \mathcal{M}}{\partial \mathbf{x}_f}
    + \left(\frac{\partial \boldsymbol{\chi}}{\partial \mathbf{x}_f}\right)^{\!T}\!\boldsymbol{\mu}.
    \label{eq:transversality_general}
\end{equation}
Each terminal state falls into one of three categories~\cite[Sec.~2.3.1]{zavoli2013}. The \emph{constrained} states $y, v_x, v_y$ leave their costates $\lambda_y(t_f), \lambda_{v_x}(t_f), \lambda_{v_y}(t_f)$ undetermined: these absorb the unknown multipliers $\boldsymbol{\mu}$ and are carried as shooting unknowns (Eq.~\ref{eq:unknowns}). The horizontal position $x$ is \emph{free and absent from the payoff}, so $\partial\mathcal{M}/\partial x = 0$ and Eq.~\eqref{eq:transversality_general} returns $\lambda_x(t_f) = 0$, consistent with $\lambda_x \equiv 0$ already obtained in Eq.~\eqref{eq:lx}. Finally, the terminal mass $m(t_f)$ is \emph{free but present in the payoff}, with $\partial\mathcal{M}/\partial m = 1$, which fixes the only nontrivial terminal costate:
\begin{equation}
    \lambda_m(t_f) = \frac{\partial \mathcal{M}}{\partial m(t_f)} = 1.
    \label{eq:trans_m}
\end{equation}
This condition carries a twofold meaning. First, $\lambda_m(t_f)$ is the sensitivity of the optimal payoff to the terminal mass (a unit increase in $m(t_f)$ improves the objective by one unit). Second, since the PMP determines $\boldsymbol{\lambda}$ only up to a positive scale factor---the Hamiltonian is homogeneous of degree one in $\boldsymbol{\lambda}$---Eq.~\eqref{eq:trans_m} fixes the normalization of the costate vector.\footnote{Here $\mathcal{M}$ is the terminal (Mayer) payoff and $\boldsymbol{\chi}$, $\boldsymbol{\mu}$ the boundary-constraint vector and its multipliers; in the general Bolza form $\mathcal{H} = \Phi + \boldsymbol{\lambda}^\top\mathbf{f}$ the symbol $\Phi$ denotes the running cost, absent in this pure-Mayer problem~\cite[Ch.~2]{zavoli2013}. The sign of the terminal costate depends on the convention used to assemble the augmented functional: the $(\mathbf{f}-\dot{\mathbf{x}})$ form adopted in Eq.~\eqref{eq:aug_payoff}, consistent with the maximization of $\mathcal{H}$, gives $\lambda_m(t_f) = +1$ and $\dot{\lambda}_m = (T/m^2)\|\boldsymbol{\lambda}_v\| \ge 0$ along a powered arc (with equality only where $T = 0$); the opposite $(\dot{\mathbf{x}}-\mathbf{f})$ form would carry a global minus sign. Both conventions yield the same physical solution; the $+1$ normalization is the one used in the shooting residual and in \texttt{ode\_burn.m}.}

Because the terminal time $t_f$ is free, the PMP supplies one further condition---the natural boundary condition for $t_f$~\cite[Ch.~2]{bryson1975}:
\begin{equation}
    \mathcal{H}(t_f) = 0.
    \label{eq:trans_tf}
\end{equation}
The system is autonomous (no explicit time dependence), so the Hamiltonian is constant along an optimal trajectory; $\mathcal{H}(t_f) = 0$ therefore implies $\mathcal{H}(t) \equiv 0$ throughout the powered arc, consistent with the optimality conditions.

\subsection{Shooting Method and Convergence}\label{sec:shooting}

Substituting the bilinear-tangent law \eqref{eq:bilinear} into Eqs.~\eqref{eq:xdot}--\eqref{eq:mdot}, the optimal control problem reduces to a \textbf{Boundary Value Problem (BVP)} solved by single shooting over the unknown costate initial conditions and the free final time.

\textbf{Two numerical refinements} keep the BVP compact and well conditioned. \emph{(i) Costate normalization.} Since the Hamiltonian is homogeneous of degree one in $\boldsymbol{\lambda}$, the costates are defined only up to a positive scale; we fix it by setting $\lambda_{m,0} = 1$, which removes $\lambda_{m,0}$ as an unknown together with the terminal condition $\lambda_m(t_f) = 1$. \emph{(ii) Hamiltonian imposed at $t_0$.} As the problem is autonomous, $\mathcal{H}$ is constant and equal to zero (free $t_f$); we therefore impose $\mathcal{H} = 0$ at the \emph{initial} instant, rather than at $t_f$ where it would accumulate integration error. Evaluating the optimal Hamiltonian~\eqref{eq:ham_opt} at $t_0$ is purely algebraic: the launch state has $v_x = v_y = 0$, so the term $\lambda_y v_y$ drops; the initial mass is $m_0 = 1$, so $T/m = T$; and with the normalization $\lambda_{m,0} = 1$ together with $Q = T/c$ at full thrust, the propellant term becomes $-\lambda_m Q = -T/c$. Collecting these contributions,
\begin{equation}
    \mathcal{H}(0) = -\lambda_{v_y,0} + T\,\lVert\boldsymbol{\lambda}_{v,0}\rVert - \frac{T}{c}
    = -\lambda_{v_y,0} + T\Bigl(\lVert\boldsymbol{\lambda}_{v,0}\rVert - \tfrac{1}{c}\Bigr) = 0,
    \label{eq:H0}
\end{equation}
where the bracket is the initial value of the switching function $S = \|\boldsymbol{\lambda}_v\|/m - \lambda_m/c$ (Section~\ref{sec:task3}), evaluated at $m_0 = 1$ and $\lambda_{m,0} = 1$. With $\lambda_x \equiv 0$ (Eq.~\ref{eq:lx}) already enforced, the shooting unknowns reduce to four,
\begin{equation}
    \boldsymbol{\zeta} = \bigl[\lambda_{v_x,0},\;\lambda_{v_y,0},\;\lambda_y,\;t_f\bigr]^T,
    \label{eq:unknowns}
\end{equation}
and $\lambda_m$ no longer enters the residual. Given a trial $\boldsymbol{\zeta}$, the state $[x, y, v_x, v_y, m]^T$ is integrated forward with \texttt{ode45} ($\mathrm{RelTol} = 10^{-10}$, $\mathrm{AbsTol} = 10^{-12}$) and the four residuals are driven to zero with \texttt{fsolve} (tolerances $10^{-10}$). The same four-unknown formulation is used for the powered arcs of Tasks~1, 2 and~4; Task~3 keeps $\lambda_m$ as a fifth unknown because the engine-cutoff (switching) condition of the coast needs it.

\textbf{Shooting function and convergence.}\ The residuals define the shooting function $\mathbf{S}:\boldsymbol{\zeta}\mapsto\mathbf{r}(\boldsymbol{\zeta})$, the mismatch of the forward-integrated trajectory against the boundary conditions, which \texttt{fsolve} drives to zero by a damped Newton iteration
\begin{equation}
    \boldsymbol{\zeta}^{k+1} = \boldsymbol{\zeta}^{k}
    - \gamma_k\left[\frac{\partial\mathbf{S}}{\partial\boldsymbol{\zeta}}\right]^{-1}\!\mathbf{S}(\boldsymbol{\zeta}^{k}),
    \label{eq:newton}
\end{equation}
with a finite-difference Jacobian (the state-transition-matrix propagation being the analytic alternative)~\cite[Sec.~2.3.2]{zavoli2013}. Single shooting is delicate precisely because the unknowns are costates with no physical meaning: a small change in $\boldsymbol{\lambda}_0$ is amplified by the forward integration into a large change in the terminal residual---the ``tail wagging the dog''---so the Jacobian is ill-conditioned, the convergence basin is small, and a cold start typically diverges. These pathologies sharpen when a switching structure is present (Task~3), where the shooting function may even become non-differentiable~\cite[Ch.~3]{zavoli2013}.

The \textbf{continuation (homotopy) strategy} is the standard remedy: the BVP is first solved at a favorable operating point ($Q \approx 3$), and each subsequent value of $Q$ is initialized with the converged solution from the neighboring point, sweeping both forward and backward across the range $Q \in [1.8, 7]$. Embedding the family of problems in the scalar parameter $Q$ supplies, at every step, an initial guess inside the local basin, so the solution branch is tracked smoothly without manual retuning. The arc-splitting introduced in Tasks~2--4 (vertical climb, coast, staging) plays an analogous role, decomposing a stiff single shoot into better-conditioned sub-arcs---the indirect-method counterpart of multiple shooting~\cite[Sec.~2.3.2]{zavoli2013}.

\end{document}
